\chapter{Introducción y metodología}
\label{chap:metodologia}

\section{Objetivo y alcance}
La actualidad en ciberseguridad genera información
abundante, pero desigual en precisión y profundidad.
Comprender un incidente requiere distinguir el mecanismo
técnico, sus consecuencias y la evidencia que respalda
cada afirmación.

Este trabajo analiza cinco casos con actividad
documentada desde el 1 de septiembre de 2026:
PHP2Shell, Ivanti Connect Secure, los servicios
gubernamentales de Moldavia, la refrigeración comercial
en Francia y las entidades financieras surcoreanas.
El objetivo es responder a la siguiente pregunta:
¿qué debilidades, consecuencias y prioridades defensivas
se identifican al comparar estos incidentes?

Se examinan su desarrollo, impacto, mitigación y
relevancia, adaptando el análisis a la información
disponible. La investigación se cierra el
8 de noviembre de 2026. La fecha del incidente
se distingue de la publicación de noticias y avisos;
las vulnerabilidades anteriores se utilizan únicamente
como antecedentes de los ataques estudiados.